\chapter{Methodik}

%%%%%

\section{Untersuchungsdesign}

%%%%%

\section{Eigener KI-Video-Datensatz}

Auswahl der Generatoren

Vorgehen der Erstellung von Testvideos

Text-Prompts und Bildvorlagen

Link auf Datensatz

Generative KI-Modelle können bereits jetzt fotorealistische Videos anhand von Text oder Deepfakes mittels Bildmaterial generieren und überwiegend im MP4-Containerformat ausgeben (siehe Table 2)

z. B. Imagine 0.9 von xAI das über die Plattform Grok angeboten wird


\begin{table}[htbp]
	\centering
	\caption{Auszug gängiger KI-Video-Generatoren zum 25.01.2026 (Eigene Darstellung)}
	\label{tab:ki_video_generatoren}
	\begin{threeparttable}
		% 4 Spalten: Modell, Anbieter und Plattform sind flexibel (X), 
		% die letzte Spalte ist zentriert (c)
		\begin{tabularx}{\textwidth}{>{\raggedright\arraybackslash}X >{\raggedright\arraybackslash}X >{\raggedright\arraybackslash}X c}
			\toprule
			\textbf{Modell (Version)} & \textbf{Anbieter} & \textbf{Plattform(en)} & \textbf{Text-/Image-to-Video} \\
			\midrule
			Veo (3.1) & Google DeepMind & Gemini, Perplexity & ja / ja \\
			Sora (2.0) & OpenAI & ChatGPT, Bing & ja / ja \\
			Pika (2.5) & Pika Labs & Pika Art & ja / ja \\
			Runway (4.5) & Runway AI & Runway & ja / ja \\
			Kling AI Video (2.6) & Kuaishou & Klingai & ja / ja \\
			Luma Dream Machine Ray (3.0) & Luma AI & Luma AI & ja / ja \\
			Imagine (0.9) & xAI & Grok & ja / ja \\
			\bottomrule
		\end{tabularx}
	\end{threeparttable}
\end{table}

%%%%%

\section{Feature-Engineering für MP4}

Strategie, um MP4-Baumstruktur abzubilden

Irrelevante Inhaltsabhängige Felder (mdat, creationTime, duration, freeSpace) rausfiltern


\section{Implementierung des Prototyps}

Inspiration durch Github-Projekt aus MSSH (Klier, 2025)\cite{klier_github_2025}
Tests mit Parsern („Online Mp4 Parser“, o. J.)
Ähnlichkeitsmetrik
